\begin{problem}[第32届全国中学生物理竞赛决赛]{55}
激光瞄准系统的设计需考虑空气折射率的变化。由于受到地表状况、海拔高度、气温、湿度和空气密度等多种因素的影响，空气的折射率在大气层中的分布是不均匀的，因而激光的传播路径并不是直线。为简化起见，假设某地的空气折射率随高度 y 的变化如下式所示
\begin{equation}
n^{2} = n_{0}^{2} + \alpha^{2} y,
\end{equation}
式中 $n_{0}$ 是 y=0 处（地面）空气的折射率，$n_{0}$ 和 $\alpha$ 均为大于零的已知常量。激光本身的传播时间可忽略。激光发射器位于坐标原点 O，如图。
\begin{subproblem}
若激光的出射方向与竖直方向 y 轴的夹角为 $\theta_{0}$（$0 \leq \theta_{0} \leq \pi/2$），求描述该激光传播路径的方程。
\end{subproblem}
\begin{subproblem}
假定目标A位于第I象限。当目标A的高度为 $y_{a}$ 时，求激光发射器可照射到的目标A的最大x-坐标值 $x_{a\max}$。
\end{subproblem}
\begin{subproblem}
激光毁伤目标需要一定的照射时间。若目标A处在激光发射器的可攻击范围内，其初始位置为 $(x_{0},y_{0})$，该目标在同一高度上以匀速度v接近激光发射器。为了使激光能始终照射该目标，激光出射角 $\theta_{0}$ 应如何随时间t而变化？
\end{subproblem}
\begin{subproblem}
激光发射器的攻击通常遵从安全击毁的原则，即既要击毁目标飞行器A，又必须使目标飞行器A水平投出的所有炸弹，都不能炸到激光发射器（炸弹在投出时相对于A静止）。假定A一旦进入激光发射器可攻击范围，激光发射器便立即用激光照射它。已知水平飞行目标 A 的高度为 $y_{a}$，击毁 A 需要激光持续照射的时间为 $t_{a}$，且位于坐标原点 O 的激光器能安全击毁它；试求 A 的速度范围。不考虑空气阻力，重力加速度大小为 g。
\end{subproblem}
\end{problem}